% GENERATED FILE. Edit canonical lessons, then run npm run generate:books.
% canonical-source-hash: fnv1a64:9c98dc9c38c5d476
% canonical-lessons: ZH-C41-duoshao, ZH-W41-duo, ZH-W41-shao, ZH-C41-qian, ZH-C41-kuai, ZH-C41-tai, ZH-C41-yao

\chapter{Paying}
\label{ch:zh-paying}
\hlchaptermodality{41}

\begin{chapteropening}
\textbf{By the end of this chapter:} \emph{I can ask how much something costs and say how much I want.}

Complete the last lesson of chapter 41: I can ask how much something costs and say how much I want.
\end{chapteropening}

\section[\texorpdfstring{\zh{多少} (duōshao)}{多少 (duōshao)}]{\zh{多少} (duōshao) — how much; how many}

\label{lesson:ZH-C41-duoshao}



\begin{quote}
Before the new one: say the Mandarin for bread, then the Mandarin for egg.
\end{quote}

\subsection*{You'll want to know: \zh{多少}}
\textbf{\zh{多少}} — \emph{duōshao} — \textquotedblleft{}how much; how many\textquotedblright{}.

\begin{quote}
\textbf{\zh{你买了多少}} \emph{nǐ mǎi le duōshao} — How many did you buy?
\end{quote}

The new character, in pieces: \zh{多} — one three-stroke piece twice, one above the other; \zh{少} — \zh{小}, the small piece, without its hook, and one long sweep below it. The next lesson puts it in your hand.

\begin{cousinweb}
\textbf{\zh{多少}} (\emph{duōshao}) is built from \zh{多} 'many' and \zh{少} 'few', both inherited from Middle Chinese, \emph{ta} and \emph{syewX} (Guangyun, in Baxter's transcription): many-few, asked as a question.
\end{cousinweb}

\begin{grammarlens}[title={What you've built}]
One more word for asking politely.
\end{grammarlens}

\subsection*{Your turn}
\begin{itemize}
\raggedright
  \item \emph{Say it:} \emph{duōshao}
  \item \emph{Say it:} \emph{duōshao}, once more
  \item \emph{Say it:} \emph{nǐ mǎi le duōshao} — How many did you buy?
  \item \emph{Recall:} say the Mandarin for bread, then the Mandarin for egg, then say \emph{duōshao} again
\end{itemize}

\begin{tcolorbox}[breakable,colback=teal!4,colframe=teal!35!black,title={Before you move on}]
What does \zh{多少} mean? (\textquotedblleft{}How much; how many\textquotedblright{}.) Say it once more.
\end{tcolorbox}
\section[\texorpdfstring{\zh{多} (duō)}{多 (duō)}]{\zh{多} — the same three strokes, twice}

\label{lesson:ZH-W41-duo}



\begin{quote}
The word first, then the character inside it.

\begin{itemize}
\raggedright
  \item \emph{Recall:} say \emph{duōshao} — \textbf{R1}, one lesson back, before it decays
\end{itemize}

\begin{quote}
\zh{了} \zh{买}
\end{quote}

\emph{Read it:} each one aloud, with nothing said first — \emph{le}, \emph{mǎi}
\end{quote}

\subsection*{Script — the same three strokes, twice}
\hlblockfigure{\detokenize{figures/ZH-W41-duo-filmstrip.pdf}}{How \zh{多} is written, stroke by stroke}

\begin{quote}
\zh{多}
\end{quote}

\textbf{One shape, twice: top, then bottom.}

\textbf{The shape, three strokes.} A short left-falling stroke. A horizontal that turns and falls to the lower left. A dot inside the turn.

\textbf{Make it again below}, a little larger and a little to the right, so the two lean together.

3 + 3 = six.

\subsection*{Your turn}
\begin{itemize}
\raggedright
  \item \emph{Write it:} \zh{多}, all six, with the model in front of you
  \item \emph{Cover:} write it again from memory, and count the pieces, not the strokes
\end{itemize}

\begin{tcolorbox}[breakable,colback=teal!4,colframe=teal!35!black,title={Before you move on}]
How many strokes in one half of \zh{多}? (\textbf{Three}.) Which half is larger? (\textbf{The lower one}.) Where does the dot go? (\textbf{Inside the turn}.)

Source: \href{https://raw.githubusercontent.com/chanind/hanzi-writer-data/68d10a4b21150cae5e1ebbd223eed289cf32d90c/data/\%E5\%A4\%9A.json}{Hanzi Writer Data \zh{多}, pinned stroke paths and medians}.
\end{tcolorbox}
\section[\texorpdfstring{\zh{少} (shǎo)}{少 (shǎo)}]{\zh{少} — \zh{小} without its hook, and a sweep}

\label{lesson:ZH-W41-shao}



\begin{quote}
The word first, then the character inside it.

\begin{itemize}
\raggedright
  \item \emph{Recall:} say \emph{duōshao} — \textbf{R1}, one lesson back, before it decays
\end{itemize}

\begin{quote}
\zh{小} \zh{买}
\end{quote}

\emph{Read it:} each one aloud, with nothing said first — \emph{xiǎo}, \emph{mǎi}
\end{quote}

\subsection*{Script — \zh{小} without its hook, and a sweep}
\hlblockfigure{\detokenize{figures/ZH-W41-shao-filmstrip.pdf}}{How \zh{少} is written, stroke by stroke}

\begin{quote}
\zh{少}
\end{quote}

\textbf{One piece, built on \zh{小}.}

\textbf{Three strokes like \zh{小}.} A vertical down the middle, but this time with no hook at the foot. A short stroke falling to the left. A dot on the right.

\textbf{One more stroke.} A long left-falling sweep from the right, under the vertical.

3 + 1 = four.

\subsection*{Your turn}
\begin{itemize}
\raggedright
  \item \emph{Write it:} \zh{少}, all four, with the model in front of you
  \item \emph{Cover:} write it again from memory, and count the pieces, not the strokes
\end{itemize}

\begin{tcolorbox}[breakable,colback=teal!4,colframe=teal!35!black,title={Before you move on}]
Which character you know is inside \zh{少}? (\textbf{\zh{小}}.) What does the middle vertical lose? (\textbf{Its hook}.) Which stroke is last? (\textbf{The long sweep to the left}.)

Source: \href{https://raw.githubusercontent.com/chanind/hanzi-writer-data/68d10a4b21150cae5e1ebbd223eed289cf32d90c/data/\%E5\%B0\%91.json}{Hanzi Writer Data \zh{少}, pinned stroke paths and medians}.
\end{tcolorbox}
\section[qián]{qián — money}

\label{lesson:ZH-C41-qian}



\begin{quote}
Before the new one: say the Mandarin for egg, then the Mandarin for how much.
\end{quote}

\subsection*{You'll want to know: qián}
\emph{qián} — \textquotedblleft{}money\textquotedblright{}.

\begin{quote}
\emph{zhège duōshao qián} — How much is this one?
\end{quote}

Its characters come later; for now it is a word for the ear.

\begin{cousinweb}
\emph{qián} is inherited from Middle Chinese \emph{dzjen} (Guangyun, in Baxter's transcription), glossed 'coin'.
\end{cousinweb}

\begin{grammarlens}[title={What you've built}]
One more word for asking politely.
\end{grammarlens}

\subsection*{Your turn}
\begin{itemize}
\raggedright
  \item \emph{Say it:} \emph{qián}
  \item \emph{Say it:} \emph{qián}, once more
  \item \emph{Say it:} \emph{zhège duōshao qián} — How much is this one?
  \item \emph{Recall:} say the Mandarin for egg, then the Mandarin for how much, then say \emph{qián} again
\end{itemize}

\begin{tcolorbox}[breakable,colback=teal!4,colframe=teal!35!black,title={Before you move on}]
What does qián mean? (\textquotedblleft{}Money\textquotedblright{}.) Say it once more.
\end{tcolorbox}
\section[kuài]{kuài — yuan; lump}

\label{lesson:ZH-C41-kuai}



\begin{quote}
Before the new one: say the Mandarin for how much, then the Mandarin for money.
\end{quote}

\subsection*{You'll want to know: kuài}
\emph{kuài} — \textquotedblleft{}yuan; lump\textquotedblright{}.

\begin{quote}
\emph{nàge shí kuài qián} — That one is ten yuan.
\end{quote}

Its characters come later; for now it is a word for the ear.

\emph{kuài} is the spoken word for the yuan: a price is a number, then \emph{kuài}, then \emph{qián}, which can be left off.

\begin{cousinweb}
\emph{kuài} is inherited from Middle Chinese \emph{khwojH} (Guangyun, in Baxter's transcription), glossed 'a clod of earth'. Its use for counting money is not given there, so how it took that job is not settled.
\end{cousinweb}

\begin{grammarlens}[title={What you've built}]
One more word for asking politely.
\end{grammarlens}

\subsection*{Your turn}
\begin{itemize}
\raggedright
  \item \emph{Say it:} \emph{kuài}
  \item \emph{Say it:} \emph{kuài}, once more
  \item \emph{Say it:} \emph{nàge shí kuài qián} — That one is ten yuan.
  \item \emph{Recall:} say the Mandarin for how much, then the Mandarin for money, then say \emph{kuài} again
\end{itemize}

\begin{tcolorbox}[breakable,colback=teal!4,colframe=teal!35!black,title={Before you move on}]
What does kuài mean? (\textquotedblleft{}Yuan; lump\textquotedblright{}.) Say it once more.
\end{tcolorbox}
\section[tài]{tài — too; too much}

\label{lesson:ZH-C41-tai}



\begin{quote}
Before the new one: say the Mandarin for money, then the Mandarin for yuan.
\end{quote}

\subsection*{You'll want to know: tài}
\emph{tài} — \textquotedblleft{}too; too much\textquotedblright{}.

\begin{quote}
\emph{tài xiǎo le} — That's too small.
\end{quote}

Its characters come later; for now it is a word for the ear.

\emph{tài} usually brings a \emph{le} at the end of the sentence with it: \emph{tài xiǎo le}.

\begin{cousinweb}
\emph{tài} is inherited from Middle Chinese \emph{thajH} (Guangyun, in Baxter's transcription), glossed 'extreme' and 'great'.
\end{cousinweb}

\begin{grammarlens}[title={What you've built}]
One more word for asking politely.
\end{grammarlens}

\subsection*{Your turn}
\begin{itemize}
\raggedright
  \item \emph{Say it:} \emph{tài}
  \item \emph{Say it:} \emph{tài}, once more
  \item \emph{Say it:} \emph{tài xiǎo le} — That's too small.
  \item \emph{Recall:} say the Mandarin for money, then the Mandarin for yuan, then say \emph{tài} again
\end{itemize}

\begin{tcolorbox}[breakable,colback=teal!4,colframe=teal!35!black,title={Before you move on}]
What does tài mean? (\textquotedblleft{}Too; too much\textquotedblright{}.) Say it once more.
\end{tcolorbox}
\section[yào]{yào — to want; to ask for}

\label{lesson:ZH-C41-yao}



\begin{quote}
Before the new one: say the Mandarin for yuan, then the Mandarin for too.
\end{quote}

\subsection*{You'll want to know: yào}
\emph{yào} — \textquotedblleft{}to want; to ask for\textquotedblright{}.

\begin{quote}
\emph{nǐ yào duōshao} — How many would you like?
\end{quote}

Its characters come later; for now it is a word for the ear.

\begin{cousinweb}
\emph{yào} is inherited from Middle Chinese \emph{'jiewH} (Guangyun, in Baxter's transcription), glossed 'to bind by agreement'. The sense 'to want' is not given there, so how it got it is not settled.
\end{cousinweb}

\begin{grammarlens}[title={What you've built}]
One more word for asking politely.
\end{grammarlens}

\subsection*{Your turn}
\begin{itemize}
\raggedright
  \item \emph{Say it:} \emph{yào}
  \item \emph{Say it:} \emph{yào}, once more
  \item \emph{Say it:} \emph{nǐ yào duōshao} — How many would you like?
  \item \emph{Recall:} say the Mandarin for yuan, then the Mandarin for too, then say \emph{yào} again
\end{itemize}

\begin{tcolorbox}[breakable,colback=teal!4,colframe=teal!35!black,title={Before you move on}]
What does yào mean? (\textquotedblleft{}To want; to ask for\textquotedblright{}.) Say it once more.
\end{tcolorbox}
